\startmode[single]
\startchapter[title={2024-11-18 - WuN}]
\stopmode

\startmode[multiple]
\startsection[title={2024-11-18}]
\stopmode

\startframedtext[location=middle,width=\textwidth]
{\bf Ablauf}

\startitemize[packed,m]
\item Tempolimit, einschränkung der Freiheit
\item Einschränkung für Arbeitende in der Stadt?
\item Tempolimit auf der Autobahn
\item Freiheit als Wahl: Jean-Paul Sarte: Tote ohne Begräbnis
\startitemize[packed]
\item Auf Teams:\\ {\em Kursmaterialien/Texte,Bilder,Links/Sartre - Tote ohne Begräbnis.pdf}
\stopitemize
\stopitemize
\stopframedtext

\startsectionlevel[title={Einschränkung für Arbeitende in der Stadt}]

\startplacetable[location={here},title={Vergleich}]
\setupTABLE[r][1][background=color,backgroundcolor=gray]
\bTABLE
\startCSV
Pro, Contra
\stopCSV
\bTR
\bTD
% Pro
\startitemize[packed]
\item Arbeitende kommen nur erschwert in die Stadt, zum Arbeitsplatz
\item Parkgebühren / Busgebühren
\item Aeltere Menschen
\item ...
\stopitemize
\eTD
\bTD
% Contra
\startitemize[packed]
\item Ausnahmen auch für den ÖPNV
\item Es handelt sich nur um die Innenstadt, autos können außerhalb parken
\item ...
\stopitemize
\eTD
\eTR
\eTABLE
\stopplacetable

\stopsectionlevel
\startsectionlevel[title={Freiheit als Wahl}]

Besprochene Zeilen
\startitemize[packed]
\item 52 - 55: Wieso sind beide "Abschaum"
\item 
\stopitemize

\stopsectionlevel
\startmode[single]
\stopchapter
\stopmode

\startmode[multiple]
\stopsection
\stopmode

\stopchapter
